\documentclass[12pt,a4paper]{article}
\usepackage[UTF8,heading=true]{ctex}
\usepackage{geometry}
\usepackage{amsmath,amssymb,amsthm}
\usepackage{graphicx}
\usepackage{booktabs}
\usepackage{longtable}
\usepackage{array}
\usepackage{enumitem}
\usepackage{hyperref}
\usepackage{fancyhdr}
\usepackage{xcolor}
\usepackage{listings}
\usepackage{setspace}

\ctexset{
  section = {
    name = {,、},
    number = \chinese{section},
    format = {\Large\bfseries\raggedright}
  }
}

\geometry{left=2.5cm,right=2.5cm,top=2.5cm,bottom=2.5cm}
\setlength{\headheight}{15pt}
\setstretch{1.5}

\pagestyle{fancy}
\fancyhf{}
\fancyhead[C]{发明专利技术交底书}
\fancyfoot[C]{\thepage}

\title{\textbf{发明专利技术交底书} \\[0.5em]
\Large 一种面向多路径推理的张量状态写时复制管理系统、芯片及方法}
\author{中国移动通信有限公司研究院}
\date{\today}

\begin{document}

\begin{center}
    {\Large\heiti 中国移动专利申请} \\[1em]
    {\Huge\heiti 技术交底书} \\[1em]
\end{center}

\noindent
\begin{tabular}{|p{3cm}|p{12cm}|}
\hline
公司编号 &  \\
\hline
发明名称 & 一种面向多路径推理的张量状态写时复制管理系统、芯片及方法 \\
\hline
申报单位 & 研究院 \\
\hline
申报类型 & 发明 \\
\hline
发明人 & 许达 \\
\hline
技术联系人 & 许达 (xudayj@chinamobile.com, +86-13521894156) \\
\hline
注意事项 & 1．技术联系人应为深入了解本申请提案技术方案的技术人员。 \\
         & 2．请按照集团公司提供的本技术交底书模板逐项填写。 \\
\hline
\end{tabular}

\vspace{1cm}

%==============================================================================
\section{发明名称}
%==============================================================================

一种面向多路径推理的张量状态写时复制管理系统、芯片及方法


%==============================================================================
\section{技术领域}
%==============================================================================

本发明涉及人工智能芯片、张量存储管理、推理加速器微架构和编译运行时协同技术领域，特别涉及一种用于树状推理、多分支搜索推理和多假设并行验证场景的张量状态写时复制管理系统、芯片及方法。

\textbf{与量子计算的关联说明：}

随着大语言模型、规划模型和具备慢思考能力的推理模型发展，模型执行形态正由单路径自回归生成转向多路径展开、回溯剪枝和分支合流。在该过程中，决定系统吞吐和成本的关键瓶颈不再只是乘加算力，而是键值缓存（KV Cache）、中间激活及状态张量的复制、回收和跨分支一致性维护。现有 GPU 与通用 AI 加速器缺少面向分支张量状态的硬件原生虚拟化机制，只能依赖软件复制和软件回收，导致显存占用、带宽开销和调度延迟急剧上升。本发明将操作系统中的写时复制思想下沉到 AI 芯片的张量状态层，属于典型的算法硅化软硬协同设计。

\textbf{关联项目：}\underline{\hspace{8cm}}（待填写）

%==============================================================================
\section{现有技术的技术方案}
%==============================================================================

\subsection{量子计算网络安全背景}

当前大模型推理已从单纯的逐 Token 生成扩展到自一致性采样、树搜索、多候选推演、验证器反馈和回溯修正等复杂流程。模型在求解数学、代码、规划和工具调用类问题时，常常需要同时维护多个假设分支，每个分支共享相同前缀上下文，但在后续若干步生成中逐步分化。

\subsection{后量子密码技术现状}

现有技术通常通过软件框架管理 KV Cache 和分支状态。常见方案包括：通过分页式注意力或块表降低连续内存要求；通过运行时在分支创建时复制指针或复制张量块；以及通过软件垃圾回收和引用计数回收被剪枝分支所占用的内存。这类方案虽然能够在软件层实现多分支推理，但其核心内存语义仍构建在“先申请、后复制、再比对、最后释放”的软件路径上，缺少直接由芯片内存控制器和指令集支持的分支状态管理机制。

%==============================================================================
\section{现有技术的缺点及本申请提案要解决的技术问题}
%==============================================================================

现有技术主要存在以下缺陷和技术问题：

(1) \textbf{分支复制成本高}：当推理框架创建新分支时，现有 GPU 往往需要复制大量 KV Cache 页、页表或块映射元数据；一旦分支数提升，复制流量会与分支数近似线性增长，导致显存带宽被状态管理而非模型计算占用。

(2) \textbf{共享前缀难以硬件级复用}：在大多数树搜索推理任务中，不同分支共享极长的前缀上下文。现有方案虽然可在软件中维护共享块，但硬件访存路径并不知道哪些张量页处于共享态，因而无法在读、写、回收三个阶段做针对性优化。

(3) \textbf{剪枝回收延迟大}：当某一逻辑分支被算法证伪时，现有技术通常需要运行时更新引用计数、回收表项、合并碎片，容易引入跨流同步和主机端参与，造成显著的尾延迟。

(4) \textbf{编译器与硬件语义割裂}：现有 ISA 中没有用于显式表达“创建逻辑分支”“终止逻辑分支”“提交分支结果”的指令，导致编译器无法把推理树结构直接映射为硬件状态机，软件框架只能以通用内存读写拼装复杂控制语义。

本申请提案旨在解决上述问题，提出一种面向多路径推理的硬件原生张量写时复制管理系统，使推理框架在创建、扩展、剪枝和合流分支时均可借助芯片级 Branch State Controller 和专用指令完成零拷贝分支创建、按写分配物理页和单周期回收。

%==============================================================================
\section{本申请提案的技术方案的详细阐述}
%==============================================================================

本发明提供一种面向多路径推理的张量状态写时复制管理系统、芯片及方法，适用于 Transformer 类模型、自回归生成模型和其他需要维护多条上下文状态链的神经网络推理平台。

\subsection{系统架构}

本发明的系统架构包括：

\textbf{分支状态管理器（Branch State Controller）}：设置于芯片内存子系统或统一地址转换子系统中，用于维护逻辑分支标识、父子关系、共享页引用计数、写脏位图和回收状态机。

\textbf{张量状态页表（Tensor State Page Table）}：用于记录 KV Cache、注意力中间状态、草稿分支临时张量和验证分支张量的逻辑页到物理页映射，支持父分支继承与子分支局部重映射。

\textbf{专用指令接口}：向编译器或运行时暴露 \texttt{FORK\_KV}、\texttt{WRITE\_KV}、\texttt{PRUNE\_KV}、\texttt{MERGE\_KV}、\texttt{QUERY\_BR} 等指令或微操作，使推理树事件能够直接转换为硬件状态跃迁。

\textbf{页粒度写时复制分配器}：当某分支首次写入共享张量页时，硬件自动分配新物理页并完成增量写入，同时保证其他分支继续读取原共享页。

\textbf{后台回收引擎}：在接收到分支剪枝或分支提交事件后，对共享计数归零的物理页进行无主机参与回收，并支持碎片整理、空闲页合并或返回至本地显存池。

\subsection{核心方法步骤}

本方案的核心流程包括以下步骤：

\subsubsection{1. 分布式密钥生成 (KeyGen)}
当推理算法需要从当前状态派生一条或多条候选思考路径时，运行时发出 \texttt{FORK\_KV} 指令。硬件不复制真实 KV 数据，而仅复制父分支的页表根指针、分支深度、页共享位图和访问权限元数据，形成子分支上下文。由此实现微小元数据复制替代大规模张量复制。

\subsubsection{2. 消息哈希与映射}
子分支在尚未生成新 Token 或尚未修改中间状态之前，对父分支共享页执行只读访问。访存单元通过分支标识查找张量状态页表，将多个逻辑分支映射到同一物理页，并通过共享计数保证数据一致性。

\subsubsection{3. 分布式高斯采样}
当某一子分支对共享页发生首次写入时，页粒度写时复制分配器侦测到该页处于共享态，随即分配新的物理页，将待写数据写入新页，并把该子分支对应页表项重映射到新页。其他分支保持对旧页的访问，从而实现“仅对实际发生差异的页分配新存储”。

\subsubsection{4. 基于 NTT 的线性变换}
当验证器、奖励模型或搜索策略判定某分支无效时，运行时发出 \texttt{PRUNE\_KV} 指令。Branch State Controller 根据该分支页表根快速遍历局部脏页列表和共享引用关系，对引用计数归零的物理页直接回收，对仍被其他分支引用的共享页仅减计数。该过程可设计为单周期触发、多周期后台完成，不要求主机端回收参与。

\subsubsection{5. 基于 Beaver 三元组的安全范数验证}
当某一子分支成为最优分支时，运行时发出 \texttt{MERGE\_KV} 或提交指令，使该子分支升级为主分支。系统仅切换活动页表根和分支上下文标识，无需执行整块张量拷贝，从而将分支合流开销压缩为元数据交换开销。

\subsubsection{6. 动态节点管理}
编译器可在图层面或调度层面识别树搜索、束搜索、自一致性采样和验证回溯等模式，在控制流边界插入专用指令。运行时则根据分支宽度、最大搜索深度和显存水位配置分支配额、页大小、预留页池和回收阈值，实现算法到芯片的直接映射。

%==============================================================================
\section{本申请提案的关键点和欲保护点}
%==============================================================================

本申请提案主要包含以下关键创新点和保护点：

(1) \textbf{将写时复制引入张量状态管理}：把传统操作系统虚拟内存中的写时复制语义迁移到 AI 推理张量状态层，形成用于 KV Cache 和中间激活的硬件级共享与分化机制。

(2) \textbf{面向分支推理的专用状态控制器}：以 Branch State Controller 维护父子分支图、页共享关系和回收状态机，使分支创建、扩展和剪枝成为内存子系统的一等公民。

(3) \textbf{专用指令驱动的硬件状态跃迁}：通过 \texttt{FORK\_KV}、\texttt{PRUNE\_KV} 和 \texttt{MERGE\_KV} 等指令，把推理树控制流直接表达为芯片可执行的微操作序列。

(4) \textbf{仅对差异页分配物理存储}：共享前缀页不重复写入，只有首次被写脏的页才触发物理新页分配，显著降低树搜索推理的显存占用和带宽流量。

(5) \textbf{硬件原生快速回收机制}：分支被剪枝时，硬件直接根据分支局部脏页列表与共享计数回收显存，减少软件垃圾回收和主机同步造成的尾延迟。

%==============================================================================
\section{与第三条中最接近的现有技术相比，本申请提案有何技术优点}
%==============================================================================

相比于现有的软件分页 KV 管理、运行时分支复制和通用显存回收机制，本方案具有以下技术优点：

(1) \textbf{显存占用显著下降}：对于共享长前缀、短差异尾部的多路径推理任务，物理页分配仅与实际差异部分相关，显著优于按分支复制整段 KV Cache 的方案，预计可使分支显存开销降低 60\% 至 90\%。

(2) \textbf{分支切换延迟更低}：分支创建和合流主要表现为页表根切换和少量元数据复制，不再依赖大块显存搬运，能够提升复杂推理的响应速度和分支宽度上限。

(3) \textbf{软件栈改动小}：无需改造矩阵乘加核心阵列，仅需在内存控制器、地址转换单元和指令解码层加入分支状态管理逻辑，工程实现路径清晰。

(4) \textbf{适配面广}：既可用于大语言模型的树搜索和多假设推理，也可用于多代理规划、候选程序搜索、投机验证树和需要保存多上下文快照的其他神经网络推理场景。

%==============================================================================
\section{其他有助于理解本申请提案的技术资料}
%==============================================================================

(1) Kwon, W., et al. ``Efficient Memory Management for Large Language Model Serving with PagedAttention.'' SOSP, 2023.
(2) Yao, S., et al. ``Tree of Thoughts: Deliberate Problem Solving with Large Language Models.'' arXiv, 2023.
(3) Wang, X., et al. ``Self-Consistency Improves Chain of Thought Reasoning in Language Models.'' ICLR, 2023.
(4) Miao, X., et al. ``SpecInfer: Accelerating Generative Large Language Model Serving with Speculative Inference and Token Tree Verification.'' arXiv, 2023.

%==============================================================================
\section{本申请提案的侵权证据可获得性}
%==============================================================================

本提案的侵权证据具有较高的可获得性：
(1) \textbf{指令与驱动接口证据}：若被诉产品对外暴露或在固件中实现与 \texttt{FORK\_KV}、\texttt{PRUNE\_KV} 类似的分支状态指令、驱动 API 或运行时接口，可作为直接证据或高相关证据。
(2) \textbf{内存行为证据}：通过性能计数器、显存带宽轨迹、分支创建时的页表变化和剪枝后的瞬时回收行为，可识别其是否采用共享页加写时复制的张量状态管理架构。
(3) \textbf{微架构与专利拆解证据}：通过芯片白皮书、驱动调试符号、固件描述、硬件性能计数器和逆向分析，可识别分支状态控制器、共享计数存储结构和后台回收状态机等专用模块。

\end{document}
